% ============================================================================
% VII. CONCLUSION (REVISED)
% ============================================================================
\section{Conclusion}
\label{sec:conclusion}

We have presented a deep learning framework for wavelength selection in multi-excitation hyperspectral imaging, addressing the challenge of identifying informative spectral layers in 4D fluorescence datasets where spatial, emission, and excitation dimensions jointly encode discriminative information.
Our approach combines three components: a 3D CAE with parallel excitation branches for unsupervised representation learning, perturbation-based attribution for tracing latent dimension importance to individual excitation-emission wavelength pairs, and Maximum Marginal Relevance selection for balancing informativeness with spectral diversity.

The framework was evaluated on three datasets covering distinct ME-HSI regimes.
On the Lichens dataset (4 classes, 192 bands), a comprehensive sweep of 3,072 configurations showed that selected wavelengths can \emph{exceed} baseline classification performance while reducing the data size: 95.2\% accuracy with 80 bands versus the 88.2\% baseline with all 192 bands, and above-baseline accuracy maintained down to 9 bands (95\% data reduction).
On the Collagen Sponges dataset (3 classes, 158 bands), the same pipeline reached 85.6\% accuracy at 30 bands (81\% reduction, 5.8 percentage points above baseline), and a ten-classifier evaluation confirmed that the selected bands transfer across classifier families.
On the Drop Data dataset (16 droplets, no labels, 214 bands), the framework was applied entirely unsupervised; the selected bands recover the three spectral types found by full-spectrum hierarchical clustering at 96.4\% per-pixel cross-validated accuracy with only 10 bands, outperforming or matching eight band-selection baselines including BS-Net-FC and a supervised Sparse-LASSO control.
Robustness analysis on Lichens confirmed genuine selection intelligence (outperforming 10,000 random combinations with $12\sigma$ separation). Object-level analysis showed that challenging specimens benefit most, with performance variance reduced by 54\%.
The 9-band optimal Lichens configuration selected one wavelength per excitation source, providing directly actionable sensor design guidelines.

To our knowledge, this work represents the first application of latent space perturbation analysis for wavelength selection in hyperspectral imaging, and the first deep learning framework specifically designed for 4D multi-excitation data.
The approach is entirely unsupervised, requiring no class labels for wavelength selection, and produces ordered rankings that can be thresholded at any desired reduction level.
The framework's modular design, separating representation learning, attribution, and selection, enables component substitution for adaptation to diverse spectroscopic applications.

As multi-excitation imaging systems become more prevalent across biomedical, agricultural, and environmental applications, principled approaches to spectral information extraction---like the framework presented here---will be essential for translating data richness into practical applications.
The demonstrated ability to improve classification accuracy while reducing the data size opens new possibilities for real-time analysis, compact sensor design, and efficient data transmission in spectral imaging applications.